% =====================================================================
%  Report — free LaTeX template by mitaa.dev
%  Compile with pdfLaTeX + Biber (Overleaf does this for you).
%  Locally:  latexmk -pdf main.tex
% =====================================================================
\documentclass[11pt, a4paper]{article}
% =====================================================================
%  Styling of the report template — mitaa.dev
%  The content lives in main.tex. Change the font, colour and
%  margins here; a regular report needs none of it.
% =====================================================================

% --- Language and font ----------------------------------------------
\usepackage[T1]{fontenc}
\usepackage[english]{babel}
\usepackage{libertinus}          % Libertinus: serif, sans and maths
\usepackage{microtype}

% --- Table of contents (load before parskip) ------------------------------
\usepackage[titles]{tocloft}  % title set as an ordinary heading
\renewcommand{\cftsecfont}{\sffamily\bfseries}
\renewcommand{\cftsecpagefont}{\sffamily\bfseries}
\setlength{\cftbeforesecskip}{0.6em}

% --- Page ------------------------------------------------------------
\usepackage[a4paper, margin=2.5cm]{geometry}
\usepackage{parskip}             % space between paragraphs, no indent
\linespread{1.08}

% --- Colour -----------------------------------------------------------
\usepackage{xcolor}
\definecolor{accent}{HTML}{1F5FA8}   % the one accent colour; change to taste
\definecolor{subtle}{HTML}{6B7280}

% --- Content ----------------------------------------------------------
\usepackage{amsmath}
\usepackage{graphicx}
\graphicspath{{figures/}}        % put your images in the figures/ folder
\usepackage{booktabs}
\usepackage{tabularx}
\usepackage{enumitem}
\setlist{itemsep=0.2em, topsep=0.3em}
\usepackage{caption}
\DeclareCaptionFont{accent}{\color{accent}}
\captionsetup{font=small, labelfont={bf,accent}, labelsep=period}
\usepackage{siunitx}             % numbers and units: \qty{9.81}{\metre\per\second\squared}

% --- References (APA 7) -------------------------------------------------
\usepackage{csquotes}
\usepackage[style=apa, backend=biber]{biblatex}
\addbibresource{references.bib}

% --- Headings ----------------------------------------------------------
\usepackage{titlesec}
\titleformat{\section}{\sffamily\Large\bfseries\color{accent}}{\thesection}{0.75em}{}
\titleformat{\subsection}{\sffamily\large\bfseries}{\thesubsection}{0.75em}{}
\titleformat{\subsubsection}{\sffamily\normalsize\bfseries}{\thesubsubsection}{0.75em}{}
\titlespacing*{\section}{0pt}{1.8em}{0.6em}
\titlespacing*{\subsection}{0pt}{1.3em}{0.3em}
\titlespacing*{\subsubsection}{0pt}{1em}{0.2em}

% --- Report details ---------------------------------------------------
% Fill these in at the top of main.tex; they appear on the title page.
% \docfield{name} creates \docname{...} to fill in and \thedocname to print it.
\newcommand{\docfield}[1]{%
  \expandafter\newcommand\csname thedoc#1\endcsname{}%
  \expandafter\newcommand\csname doc#1\endcsname[1]{%
    \expandafter\renewcommand\csname thedoc#1\endcsname{##1}}}
\docfield{title}
\docfield{subtitle}
\docfield{authors}
\docfield{programme}
\docfield{course}
\docfield{supervisor}
\docfield{institution}
\docfield{version}
\docfield{date}
\docdate{\today}

% --- Header and footer ------------------------------------------------
\usepackage{fancyhdr}
\pagestyle{fancy}
\fancyhf{}
\fancyhead[L]{\small\sffamily\color{subtle}\thedoctitle}
\fancyhead[R]{\small\sffamily\color{subtle}\nouppercase{\leftmark}}
\fancyfoot[R]{\small\sffamily\thepage}
\renewcommand{\headrulewidth}{0pt}
\renewcommand{\sectionmark}[1]{\markboth{#1}{}}
\fancypagestyle{plain}{\fancyhf{}\fancyfoot[R]{\small\sffamily\thepage}}

% --- Title page -----------------------------------------------------------
% One line of details; empty fields are skipped.
\newcommand{\docmeta}[2]{%
  \ifdefempty{#2}{}{%
    \noindent\makebox[8em][l]{\color{subtle}#1}%
    \begin{tabular}[t]{@{}l@{}}#2\end{tabular}\par}}

\newcommand{\maketitlepage}{%
  \begin{titlepage}
    \raggedright\sffamily
    {\color{accent}\rule{3em}{3pt}}\par
    \vspace{2.5em}
    {\small\bfseries\color{accent}REPORT\par}
    \vspace{0.8em}
    {\fontsize{30}{35}\selectfont\bfseries\thedoctitle\par}
    \ifdefempty{\thedocsubtitle}{}{%
      \vspace{0.8em}{\Large\color{subtle}\thedocsubtitle\par}}
    \vfill
    {\setlength{\parskip}{0.35em}%
      \docmeta{Authors}{\thedocauthors}
      \docmeta{Programme}{\thedocprogramme}
      \docmeta{Course}{\thedoccourse}
      \docmeta{Lecturer}{\thedocsupervisor}
      \docmeta{Institution}{\thedocinstitution}
      \docmeta{Version}{\thedocversion}
      \docmeta{Date}{\thedocdate}}
  \end{titlepage}}

% --- Links (load last) -----------------------------------------
\usepackage[hidelinks]{hyperref}
\usepackage[nameinlink, noabbrev]{cleveref}


% --- Details: fill these in -------------------------------------------
\doctitle{The Title of Your Report}
\docsubtitle{A subtitle that says what was studied or built}
\docauthors{First Last (1234567) \\ Second Student (7654321)}
\docprogramme{Degree programme}
\doccourse{Course or project name}
\docsupervisor{Lecturer's name}
\docinstitution{University or college}
\docversion{1.0}
\docdate{\today}

\begin{document}

\maketitlepage

% --- Front matter: Roman page numbers -------------------------------------
\pagenumbering{roman}

\section*{Summary}
Summarise the whole report in one paragraph of 150 to 250 words: the aim,
the approach, the main results and the conclusion. Write it last; anyone
who reads only the summary should know what you found.

\tableofcontents
\clearpage

% --- Main text: Arabic page numbers -----------------------------------------
\pagenumbering{arabic}

% ---------------------------------------------------------------------
\section{Introduction}
\label{sec:introduction}
% ---------------------------------------------------------------------
Describe the background: why was this report written, and for whom? Then
state the aim in one sentence and, if there is one, the question you
answer.

\subsection{Aim}
The aim of this report is \ldots{} Make the aim measurable, so that
\cref{sec:conclusion} can show whether it was met.

\subsection{Outline}
\Cref{sec:theory} gives the background, \cref{sec:method} describes the
method and \cref{sec:results} the results. These are discussed in
\cref{sec:discussion}, after which \cref{sec:conclusion} concludes.

% ---------------------------------------------------------------------
\section{Theoretical background}
\label{sec:theory}
% ---------------------------------------------------------------------
Explain the concepts and formulas the reader needs to follow the rest, and
cite your sources \parencite{lamport1994}. A numbered equation looks like
this:
\begin{equation}
  \label{eq:ohm}
  U = I \cdot R
\end{equation}
where $U$ is the voltage in volts, $I$ the current in amperes and $R$ the
resistance in ohms. In the text you refer to it as \cref{eq:ohm}.

% ---------------------------------------------------------------------
\section{Method}
\label{sec:method}
% ---------------------------------------------------------------------
Describe your work precisely enough that someone else could repeat it.

\subsection{Materials}
\begin{itemize}
  \item an adjustable power supply from \qty{0}{\volt} to \qty{30}{\volt};
  \item a multimeter with an accuracy of \qty{0.5}{\percent};
  \item resistors from \qty{100}{\ohm} to \qty{10}{\kilo\ohm}.
\end{itemize}

\subsection{Procedure}
\begin{enumerate}
  \item Build the circuit following the diagram in \cref{fig:setup}.
  \item Set the voltage and read the current.
  \item Repeat the measurement three times and record the mean.
\end{enumerate}

\begin{figure}[htbp]
  \centering
  \includegraphics[width=0.6\linewidth]{example-image-a}
  \caption{The measurement set-up. Replace this example with your own
    image in the \texttt{figures/} folder.}
  \label{fig:setup}
\end{figure}

% ---------------------------------------------------------------------
\section{Results}
\label{sec:results}
% ---------------------------------------------------------------------
Present the results without explaining them yet. \Cref{tab:measurements}
gives the measured values; the raw data are in \cref{app:data}.

\begin{table}[htbp]
  \centering
  \caption{Measured current at different voltages.}
  \label{tab:measurements}
  \begin{tabular}{S[table-format=2.1] S[table-format=3.1] S[table-format=1.2]}
    \toprule
    {Voltage (\si{\volt})} & {Current (\si{\milli\ampere})} & {Deviation (\si{\percent})} \\
    \midrule
     5.0 &  49.8 & 0.40 \\
    10.0 & 100.3 & 0.30 \\
    15.0 & 149.1 & 0.60 \\
    20.0 & 200.6 & 0.30 \\
    \bottomrule
  \end{tabular}
\end{table}

% ---------------------------------------------------------------------
\section{Discussion}
\label{sec:discussion}
% ---------------------------------------------------------------------
Explain what the results mean. Do they agree with the theory in
\cref{sec:theory}? Which sources of error play a part, and how reliable
are the measurements? Be honest about what could be better.

% ---------------------------------------------------------------------
\section{Conclusion and recommendations}
\label{sec:conclusion}
% ---------------------------------------------------------------------
Answer the aim from \cref{sec:introduction} in a few sentences, without
new information. Close with concrete recommendations: what should happen,
by whom, and why?

\printbibliography[title={References}, heading=bibintoc]

% --- Appendices -------------------------------------------------------------
\clearpage
\appendix

\section{Measurement data}
\label{app:data}
Put data here that would make the main text too long: raw measurements,
full tables or extended calculations.

\section{Calculations}
Work through an example calculation step by step, for instance with
\cref{eq:ohm}: at $U = \qty{10}{\volt}$ and $R = \qty{100}{\ohm}$,
$I = U / R = \qty{0.1}{\ampere}$.

\end{document}
